% Datapath del TRAPECIO y de su MESETA, registro por registro: mca_trap
% expandido y la cadena que convierte su salida en F_TRAP.
%   mca_trap     las tres líneas de retardo (memorias circulares con un solo
%                puntero de escritura wp), las banderas de llenado desde el clr
%                (cnt contra k, l y k+l), la doble diferencia, el primer
%                integrador, el polo-cero, el segundo integrador, el
%                corrimiento y la saturación: 8 muestras de latencia, en dos
%                filas con la regla de muestras arriba de cada una
%   mca_window   la ventana de cada forma en el marco de la ranura:
%                ini = off + 8, fin = ini + 2^lw, saturadas a 16 bits
%   mca_slot     la suma del trapecio mientras acumula y ini <= len < fin, y
%                cuántas muestras sumó
%   mca_publish  el promedio (suma >>> lw) de la forma elegida f*, saturado, y
%                cero si la ventana no se completó o el pulso quedó marcado
% Fuentes: modulos/comun/rtl/mca_trap.sv, modulos/mca/rtl/mca_window.sv,
% mca_slot.sv y mca_publish.sv (feat_trap), con los parámetros de
% rtl/mine/tops/mca_red_pitaya_top.sv (3 formas, 3 ranuras). La vista por
% bloques está en panorama_mca_window.tex y panorama_mca_slot.tex; cómo se
% elige f*, en trapecio/banco_eleccion_forma.tex.
% Estilo: docs/figuras/estilo_tikz.tex, con los símbolos de operación. El color
% del fondo es la función, el mismo de los panoramas del MCA (azul
% preprocesamiento, verde estadísticos, naranja publicación); cada fondo es
% también un módulo. Nombres de red sin cable (k, l, Mq, sh, clr, val, wp, a_k,
% a_l, fk1, fl1, fkl2, kl, ini, fin, len, acumula, abre, en, f*, lw, marcado):
% la señal del mismo nombre, que se dibuja una sola vez donde nace, o que
% viene de otra figura.
% El epígrafe no va en la figura: está en datapath_trap_meseta.txt y en el doc
% que la incluye.
%
% Compilar:   pdflatex datapath_trap_meseta.tex
%             pdftocairo -svg datapath_trap_meseta.pdf datapath_trap_meseta.svg
%
\documentclass[border=0pt,tikz]{standalone}
\usepackage[utf8]{inputenc}
\usepackage[T1]{fontenc}
\usepackage{lmodern}
\usepackage[none]{hyphenat}
% Estilo común de las figuras TikZ (panoramas y diagramas de bloques).
%
% Cada figura lo incluye con % Estilo común de las figuras TikZ (panoramas y diagramas de bloques).
%
% Cada figura lo incluye con % Estilo común de las figuras TikZ (panoramas y diagramas de bloques).
%
% Cada figura lo incluye con \input{../../figuras/estilo_tikz.tex} (las figuras
% viven en docs/<módulo>/figuras/). Un solo lugar para la paleta, los estilos de
% bloque y el formato de los parámetros: si cambia algo acá, cambia en todas.
%
% La paleta es la misma de docs/figuras/estilo_svg.py y de PALETA en
% modulos/mk/datasheet.py; si se toca un color, se toca en los tres.
%
% Requiere: \usepackage{tikz} (lo da la clase standalone con la opción tikz).

\usetikzlibrary{arrows.meta,positioning,calc,fit,backgrounds,shapes.multipart,
                decorations.markings,shadows}

% ---------------------------------------------------------------- colores ---
\definecolor{azul}{RGB}{108,172,228}      % acondicionamiento
\definecolor{verde}{RGB}{110,214,142}     % captura / histogramas en BRAM
\definecolor{amarillo}{RGB}{252,232,160}  % trigger
\definecolor{naranja}{RGB}{252,206,160}   % camino a DDR
\definecolor{grisr}{RGB}{226,226,226}     % compartido / infraestructura
\definecolor{rojo}{RGB}{240,128,128}      % bus del sistema

% ---------------------------------------------------------------- estilos ---
\tikzset{
  >={Stealth[length=2.4mm,width=1.8mm]},
  % fondo blanco siempre: el SVG se ve igual en modo claro y oscuro
  fondo/.style   = {show background rectangle, inner frame sep=8pt,
                    background rectangle/.style={fill=white, draw=none}},
  bloque/.style  = {draw=black, line width=0.6pt, rounded corners=2pt,
                    align=center, text width=2.2cm, minimum height=1.15cm,
                    inner sep=2pt, font=\scriptsize, fill=#1},
  datos/.style   = {bloque=azul!35},
  trig/.style    = {bloque=amarillo},
  captura/.style = {bloque=verde!55},
  ddr/.style     = {bloque=naranja!70},
  comp/.style    = {bloque=grisr},
  bus/.style     = {bloque=rojo!45},
  % lógica escrita dentro del módulo, sin submódulo propio
  etapa/.style   = {bloque=#1, line width=0.5pt},
  etapa/.default = verde!20,
  % lo que el bitstream no trae (atado a cero, stub, salida muda)
  mudo/.style    = {bloque=white, draw=black!45, dashed, text=black!55},
  %
  % BLOQUE CON PARÁMETROS, en el formato del datasheet: el encabezado lleva el
  % nombre y los parámetros con que lo instancia el padre (NOMBRE=valor, en
  % monoespaciada, uno por línea); el cuerpo, más claro, lo que hace.
  %   \node[datos, param={azul!35}{azul!12}] {\blk{nombre}{} \prm{A=1\\B=2}
  %                                            \nodepart{two} función};
  param/.style 2 args = {rectangle split, rectangle split parts=2,
                         rectangle split part fill={#1, #2}},
  %
  flujo/.style   = {->, line width=0.8pt},
  ctrl/.style    = {->, line width=0.6pt, black!65},
  % la línea que pasa por encima en un cruce: un halo blanco la separa
  cruce/.style   = {preaction={draw=white, line width=3.4pt, -}},
  % marca de ancho de bus: barra oblicua a mitad de la flecha, con el número
  % encima. \draw[flujo, ancho=14] ...; para otra posición, ancho en=0.3.
  ancho/.style   = {decoration={markings, mark=at position \anchopos with {
                      \draw[line width=0.6pt, -] (-0.1,-0.14) -- (0.1,0.14);
                      \node[font=\scriptsize, above=1.5pt] {#1};}},
                    postaction={decorate}},
  % la misma marca en una flecha vertical: el número va a la derecha
  anchov/.style  = {decoration={markings, mark=at position \anchopos with {
                      \draw[line width=0.6pt, -] (-0.14,-0.1) -- (0.14,0.1);
                      \node[font=\scriptsize, right=3pt] {#1};}},
                    postaction={decorate}},
  ancho en/.code = {\def\anchopos{#1}},
  sen/.style     = {font=\scriptsize\ttfamily, inner sep=1.5pt},
  nodo/.style    = {font=\tiny\ttfamily, inner sep=1pt, black!75},
  pin/.style     = {font=\scriptsize, align=center, inner sep=1pt},
  grupo/.style   = {draw=black!55, dashed, line width=0.7pt, rounded corners=4pt},
  canal/.style   = {grupo},
  contenedor/.style = {draw=black!80, line width=1.0pt, rounded corners=6pt},
  % marca numerada: cada número es una fila de la tabla del documento
  marca/.style   = {circle, fill=black, text=white, font=\scriptsize\bfseries,
                    inner sep=0pt, minimum size=4.3mm},
  %
  % DECISIONES (docs/mca/vista_general_mca.md). Cada opción es un bloque, todas
  % entran a un mux (\muxdec, abajo) y el registro que elige entra por abajo.
  % la opción que trae el reset (el comportamiento histórico): borde doble
  defecto/.style = {double=white, double distance=0.8pt, line width=0.45pt},
  % rótulo de la decisión, el mismo Dn del panorama y del documento
  decision/.style = {circle, fill=black, text=white, font=\small\bfseries,
                     inner sep=0pt, minimum size=7mm},
  %
  % LÓGICA (docs/mca/figuras/contabilidad/ y los detalles de las decisiones)
  % registro: caja gris con el triángulo de reloj en el borde de abajo
  registro/.style = {draw=black, line width=0.6pt, fill=grisr, align=center,
                     font=\scriptsize, inner sep=2pt, minimum width=#1,
                     minimum height=0.85cm,
                     path picture={\draw[line width=0.5pt]
                       ([xshift=-1.1mm]path picture bounding box.south) --
                       ([yshift=1.6mm]path picture bounding box.south) --
                       ([xshift=1.1mm]path picture bounding box.south);}},
  registro/.default = 0.8cm,
  % condición de una cadena if / else-if: "sí" sale a la derecha, "no" abajo
  pregunta/.style = {draw=black, line width=0.5pt, rounded corners=7pt,
                     fill=white, font=\scriptsize\ttfamily, align=center,
                     inner sep=2.5pt, minimum height=0.62cm, text width=#1},
  pregunta/.default = 1.9cm,
  si/.style      = {font=\tiny, inner sep=1pt, above, pos=0.12},
  no/.style      = {font=\tiny, inner sep=1pt, left, pos=0.45},
  %
  titulo/.style  = {anchor=north west, font=\small, inner sep=4pt, align=left},
  rotulo/.style  = {anchor=north west, font=\scriptsize\itshape, black!70,
                    inner sep=2pt, align=left},
}

% submódulo: nombre en monoespaciada; lógica propia: nombre en romana
\newcommand{\blk}[2]{{\ttfamily\bfseries #1}\\[1pt]#2}
\newcommand{\etq}[2]{{\bfseries #1}\\[1pt]#2}
\newcommand{\punto}[1]{\fill (#1) circle (1.3pt);}

% Parámetros en el formato del datasheet: NOMBRE=valor en monoespaciada, una
% entrada por celda, hasta cuatro columnas. Una columna: \prm{A=1\\B=2};
% dos: \prm{A=1 & B=2\\C=3}.
\newcommand{\prm}[1]{{\ttfamily\fontsize{5.6}{6.6}\selectfont
  \begin{tabular}{@{}l@{\hspace{7pt}}l@{\hspace{7pt}}l@{\hspace{7pt}}l@{}}#1\end{tabular}}}

\def\anchopos{0.5}

% Mux de una decisión: trapecio con el lado ancho (las entradas) a la izquierda.
% #1 = nombre, #2 = centro, #3 = alto en cm. Deja tres coordenadas:
%   (#1-in)  centro del lado de entradas: las flechas llegan a (#1-in |- opcion.east)
%   (#1-out) la salida
%   (#1-sel) mitad del lado de abajo: por ahí entra el registro que elige
% El rótulo de adentro es opcional: \muxdec{m}{(4,0)}{3.6} y después
% \node[...] at (m) {...}.
\newcommand{\muxdec}[3]{%
  \path #2 coordinate (#1);
  \draw[line width=0.6pt, fill=grisr]
    ($(#1)+(-0.45,{#3/2})$) -- ($(#1)+(0.45,{#3/2-0.35})$) --
    ($(#1)+(0.45,{-#3/2+0.35})$) -- ($(#1)+(-0.45,{-#3/2})$) -- cycle;
  \coordinate (#1-in)  at ($(#1)+(-0.45,0)$);
  \coordinate (#1-out) at ($(#1)+(0.45,0)$);
  \coordinate (#1-sel) at ($(#1)+(0,{-#3/2+0.175})$);}

% El mismo mux rotado 90 grados, para un flujo de arriba hacia abajo: entradas
% por el lado ancho de arriba, salida abajo y el registro que elige por la
% izquierda. #1 = nombre, #2 = centro, #3 = ancho en cm. Deja (#1-in) centro del
% lado de entradas, (#1-out) la salida y (#1-sel) mitad del lado izquierdo.
\newcommand{\muxdecv}[3]{%
  \path #2 coordinate (#1);
  \draw[line width=0.6pt, fill=grisr]
    ($(#1)+({-#3/2},0.45)$) -- ($(#1)+({#3/2},0.45)$) --
    ($(#1)+({#3/2-0.35},-0.45)$) -- ($(#1)+({-#3/2+0.35},-0.45)$) -- cycle;
  \coordinate (#1-in)  at ($(#1)+(0,0.45)$);
  \coordinate (#1-out) at ($(#1)+(0,-0.45)$);
  \coordinate (#1-sel) at ($(#1)+({-#3/2+0.175},0)$);}

% Leyenda de colores + nota. #1 = posición, #2 = color/texto, #3 = nota.
\newcommand{\leyenda}[3]{%
  \begin{scope}[shift={(#1)}, every node/.style={font=\scriptsize}]
    \coordinate (t0) at (-0.5,0);
    \foreach \c/\t [count=\i] in {#2} {
      \pgfmathtruncatemacro\j{\i-1}
      \node[draw, fill=\c, minimum width=4mm, minimum height=3mm, inner sep=0pt,
            right=5mm of t\j.east, anchor=west] (l\i) {};
      \node[right=1.5mm of l\i] (t\i) {\t};
    }
    \node[anchor=north west, black!70, text width=19cm, align=left]
      at (-0.05,-0.3) {#3};
  \end{scope}}

% Lo que dicen los parámetros, para la nota de la leyenda
\newcommand{\notaparam}{\texttt{NOMBRE=valor}: parámetros con que el padre
  instancia el bloque (se omiten los que sólo pasan el del padre sin cambio).}
 (las figuras
% viven en docs/<módulo>/figuras/). Un solo lugar para la paleta, los estilos de
% bloque y el formato de los parámetros: si cambia algo acá, cambia en todas.
%
% La paleta es la misma de docs/figuras/estilo_svg.py y de PALETA en
% modulos/mk/datasheet.py; si se toca un color, se toca en los tres.
%
% Requiere: \usepackage{tikz} (lo da la clase standalone con la opción tikz).

\usetikzlibrary{arrows.meta,positioning,calc,fit,backgrounds,shapes.multipart,
                decorations.markings,shadows}

% ---------------------------------------------------------------- colores ---
\definecolor{azul}{RGB}{108,172,228}      % acondicionamiento
\definecolor{verde}{RGB}{110,214,142}     % captura / histogramas en BRAM
\definecolor{amarillo}{RGB}{252,232,160}  % trigger
\definecolor{naranja}{RGB}{252,206,160}   % camino a DDR
\definecolor{grisr}{RGB}{226,226,226}     % compartido / infraestructura
\definecolor{rojo}{RGB}{240,128,128}      % bus del sistema

% ---------------------------------------------------------------- estilos ---
\tikzset{
  >={Stealth[length=2.4mm,width=1.8mm]},
  % fondo blanco siempre: el SVG se ve igual en modo claro y oscuro
  fondo/.style   = {show background rectangle, inner frame sep=8pt,
                    background rectangle/.style={fill=white, draw=none}},
  bloque/.style  = {draw=black, line width=0.6pt, rounded corners=2pt,
                    align=center, text width=2.2cm, minimum height=1.15cm,
                    inner sep=2pt, font=\scriptsize, fill=#1},
  datos/.style   = {bloque=azul!35},
  trig/.style    = {bloque=amarillo},
  captura/.style = {bloque=verde!55},
  ddr/.style     = {bloque=naranja!70},
  comp/.style    = {bloque=grisr},
  bus/.style     = {bloque=rojo!45},
  % lógica escrita dentro del módulo, sin submódulo propio
  etapa/.style   = {bloque=#1, line width=0.5pt},
  etapa/.default = verde!20,
  % lo que el bitstream no trae (atado a cero, stub, salida muda)
  mudo/.style    = {bloque=white, draw=black!45, dashed, text=black!55},
  %
  % BLOQUE CON PARÁMETROS, en el formato del datasheet: el encabezado lleva el
  % nombre y los parámetros con que lo instancia el padre (NOMBRE=valor, en
  % monoespaciada, uno por línea); el cuerpo, más claro, lo que hace.
  %   \node[datos, param={azul!35}{azul!12}] {\blk{nombre}{} \prm{A=1\\B=2}
  %                                            \nodepart{two} función};
  param/.style 2 args = {rectangle split, rectangle split parts=2,
                         rectangle split part fill={#1, #2}},
  %
  flujo/.style   = {->, line width=0.8pt},
  ctrl/.style    = {->, line width=0.6pt, black!65},
  % la línea que pasa por encima en un cruce: un halo blanco la separa
  cruce/.style   = {preaction={draw=white, line width=3.4pt, -}},
  % marca de ancho de bus: barra oblicua a mitad de la flecha, con el número
  % encima. \draw[flujo, ancho=14] ...; para otra posición, ancho en=0.3.
  ancho/.style   = {decoration={markings, mark=at position \anchopos with {
                      \draw[line width=0.6pt, -] (-0.1,-0.14) -- (0.1,0.14);
                      \node[font=\scriptsize, above=1.5pt] {#1};}},
                    postaction={decorate}},
  % la misma marca en una flecha vertical: el número va a la derecha
  anchov/.style  = {decoration={markings, mark=at position \anchopos with {
                      \draw[line width=0.6pt, -] (-0.14,-0.1) -- (0.14,0.1);
                      \node[font=\scriptsize, right=3pt] {#1};}},
                    postaction={decorate}},
  ancho en/.code = {\def\anchopos{#1}},
  sen/.style     = {font=\scriptsize\ttfamily, inner sep=1.5pt},
  nodo/.style    = {font=\tiny\ttfamily, inner sep=1pt, black!75},
  pin/.style     = {font=\scriptsize, align=center, inner sep=1pt},
  grupo/.style   = {draw=black!55, dashed, line width=0.7pt, rounded corners=4pt},
  canal/.style   = {grupo},
  contenedor/.style = {draw=black!80, line width=1.0pt, rounded corners=6pt},
  % marca numerada: cada número es una fila de la tabla del documento
  marca/.style   = {circle, fill=black, text=white, font=\scriptsize\bfseries,
                    inner sep=0pt, minimum size=4.3mm},
  %
  % DECISIONES (docs/mca/vista_general_mca.md). Cada opción es un bloque, todas
  % entran a un mux (\muxdec, abajo) y el registro que elige entra por abajo.
  % la opción que trae el reset (el comportamiento histórico): borde doble
  defecto/.style = {double=white, double distance=0.8pt, line width=0.45pt},
  % rótulo de la decisión, el mismo Dn del panorama y del documento
  decision/.style = {circle, fill=black, text=white, font=\small\bfseries,
                     inner sep=0pt, minimum size=7mm},
  %
  % LÓGICA (docs/mca/figuras/contabilidad/ y los detalles de las decisiones)
  % registro: caja gris con el triángulo de reloj en el borde de abajo
  registro/.style = {draw=black, line width=0.6pt, fill=grisr, align=center,
                     font=\scriptsize, inner sep=2pt, minimum width=#1,
                     minimum height=0.85cm,
                     path picture={\draw[line width=0.5pt]
                       ([xshift=-1.1mm]path picture bounding box.south) --
                       ([yshift=1.6mm]path picture bounding box.south) --
                       ([xshift=1.1mm]path picture bounding box.south);}},
  registro/.default = 0.8cm,
  % condición de una cadena if / else-if: "sí" sale a la derecha, "no" abajo
  pregunta/.style = {draw=black, line width=0.5pt, rounded corners=7pt,
                     fill=white, font=\scriptsize\ttfamily, align=center,
                     inner sep=2.5pt, minimum height=0.62cm, text width=#1},
  pregunta/.default = 1.9cm,
  si/.style      = {font=\tiny, inner sep=1pt, above, pos=0.12},
  no/.style      = {font=\tiny, inner sep=1pt, left, pos=0.45},
  %
  titulo/.style  = {anchor=north west, font=\small, inner sep=4pt, align=left},
  rotulo/.style  = {anchor=north west, font=\scriptsize\itshape, black!70,
                    inner sep=2pt, align=left},
}

% submódulo: nombre en monoespaciada; lógica propia: nombre en romana
\newcommand{\blk}[2]{{\ttfamily\bfseries #1}\\[1pt]#2}
\newcommand{\etq}[2]{{\bfseries #1}\\[1pt]#2}
\newcommand{\punto}[1]{\fill (#1) circle (1.3pt);}

% Parámetros en el formato del datasheet: NOMBRE=valor en monoespaciada, una
% entrada por celda, hasta cuatro columnas. Una columna: \prm{A=1\\B=2};
% dos: \prm{A=1 & B=2\\C=3}.
\newcommand{\prm}[1]{{\ttfamily\fontsize{5.6}{6.6}\selectfont
  \begin{tabular}{@{}l@{\hspace{7pt}}l@{\hspace{7pt}}l@{\hspace{7pt}}l@{}}#1\end{tabular}}}

\def\anchopos{0.5}

% Mux de una decisión: trapecio con el lado ancho (las entradas) a la izquierda.
% #1 = nombre, #2 = centro, #3 = alto en cm. Deja tres coordenadas:
%   (#1-in)  centro del lado de entradas: las flechas llegan a (#1-in |- opcion.east)
%   (#1-out) la salida
%   (#1-sel) mitad del lado de abajo: por ahí entra el registro que elige
% El rótulo de adentro es opcional: \muxdec{m}{(4,0)}{3.6} y después
% \node[...] at (m) {...}.
\newcommand{\muxdec}[3]{%
  \path #2 coordinate (#1);
  \draw[line width=0.6pt, fill=grisr]
    ($(#1)+(-0.45,{#3/2})$) -- ($(#1)+(0.45,{#3/2-0.35})$) --
    ($(#1)+(0.45,{-#3/2+0.35})$) -- ($(#1)+(-0.45,{-#3/2})$) -- cycle;
  \coordinate (#1-in)  at ($(#1)+(-0.45,0)$);
  \coordinate (#1-out) at ($(#1)+(0.45,0)$);
  \coordinate (#1-sel) at ($(#1)+(0,{-#3/2+0.175})$);}

% El mismo mux rotado 90 grados, para un flujo de arriba hacia abajo: entradas
% por el lado ancho de arriba, salida abajo y el registro que elige por la
% izquierda. #1 = nombre, #2 = centro, #3 = ancho en cm. Deja (#1-in) centro del
% lado de entradas, (#1-out) la salida y (#1-sel) mitad del lado izquierdo.
\newcommand{\muxdecv}[3]{%
  \path #2 coordinate (#1);
  \draw[line width=0.6pt, fill=grisr]
    ($(#1)+({-#3/2},0.45)$) -- ($(#1)+({#3/2},0.45)$) --
    ($(#1)+({#3/2-0.35},-0.45)$) -- ($(#1)+({-#3/2+0.35},-0.45)$) -- cycle;
  \coordinate (#1-in)  at ($(#1)+(0,0.45)$);
  \coordinate (#1-out) at ($(#1)+(0,-0.45)$);
  \coordinate (#1-sel) at ($(#1)+({-#3/2+0.175},0)$);}

% Leyenda de colores + nota. #1 = posición, #2 = color/texto, #3 = nota.
\newcommand{\leyenda}[3]{%
  \begin{scope}[shift={(#1)}, every node/.style={font=\scriptsize}]
    \coordinate (t0) at (-0.5,0);
    \foreach \c/\t [count=\i] in {#2} {
      \pgfmathtruncatemacro\j{\i-1}
      \node[draw, fill=\c, minimum width=4mm, minimum height=3mm, inner sep=0pt,
            right=5mm of t\j.east, anchor=west] (l\i) {};
      \node[right=1.5mm of l\i] (t\i) {\t};
    }
    \node[anchor=north west, black!70, text width=19cm, align=left]
      at (-0.05,-0.3) {#3};
  \end{scope}}

% Lo que dicen los parámetros, para la nota de la leyenda
\newcommand{\notaparam}{\texttt{NOMBRE=valor}: parámetros con que el padre
  instancia el bloque (se omiten los que sólo pasan el del padre sin cambio).}
 (las figuras
% viven en docs/<módulo>/figuras/). Un solo lugar para la paleta, los estilos de
% bloque y el formato de los parámetros: si cambia algo acá, cambia en todas.
%
% La paleta es la misma de docs/figuras/estilo_svg.py y de PALETA en
% modulos/mk/datasheet.py; si se toca un color, se toca en los tres.
%
% Requiere: \usepackage{tikz} (lo da la clase standalone con la opción tikz).

\usetikzlibrary{arrows.meta,positioning,calc,fit,backgrounds,shapes.multipart,
                decorations.markings,shadows}

% ---------------------------------------------------------------- colores ---
\definecolor{azul}{RGB}{108,172,228}      % acondicionamiento
\definecolor{verde}{RGB}{110,214,142}     % captura / histogramas en BRAM
\definecolor{amarillo}{RGB}{252,232,160}  % trigger
\definecolor{naranja}{RGB}{252,206,160}   % camino a DDR
\definecolor{grisr}{RGB}{226,226,226}     % compartido / infraestructura
\definecolor{rojo}{RGB}{240,128,128}      % bus del sistema

% ---------------------------------------------------------------- estilos ---
\tikzset{
  >={Stealth[length=2.4mm,width=1.8mm]},
  % fondo blanco siempre: el SVG se ve igual en modo claro y oscuro
  fondo/.style   = {show background rectangle, inner frame sep=8pt,
                    background rectangle/.style={fill=white, draw=none}},
  bloque/.style  = {draw=black, line width=0.6pt, rounded corners=2pt,
                    align=center, text width=2.2cm, minimum height=1.15cm,
                    inner sep=2pt, font=\scriptsize, fill=#1},
  datos/.style   = {bloque=azul!35},
  trig/.style    = {bloque=amarillo},
  captura/.style = {bloque=verde!55},
  ddr/.style     = {bloque=naranja!70},
  comp/.style    = {bloque=grisr},
  bus/.style     = {bloque=rojo!45},
  % lógica escrita dentro del módulo, sin submódulo propio
  etapa/.style   = {bloque=#1, line width=0.5pt},
  etapa/.default = verde!20,
  % lo que el bitstream no trae (atado a cero, stub, salida muda)
  mudo/.style    = {bloque=white, draw=black!45, dashed, text=black!55},
  %
  % BLOQUE CON PARÁMETROS, en el formato del datasheet: el encabezado lleva el
  % nombre y los parámetros con que lo instancia el padre (NOMBRE=valor, en
  % monoespaciada, uno por línea); el cuerpo, más claro, lo que hace.
  %   \node[datos, param={azul!35}{azul!12}] {\blk{nombre}{} \prm{A=1\\B=2}
  %                                            \nodepart{two} función};
  param/.style 2 args = {rectangle split, rectangle split parts=2,
                         rectangle split part fill={#1, #2}},
  %
  flujo/.style   = {->, line width=0.8pt},
  ctrl/.style    = {->, line width=0.6pt, black!65},
  % la línea que pasa por encima en un cruce: un halo blanco la separa
  cruce/.style   = {preaction={draw=white, line width=3.4pt, -}},
  % marca de ancho de bus: barra oblicua a mitad de la flecha, con el número
  % encima. \draw[flujo, ancho=14] ...; para otra posición, ancho en=0.3.
  ancho/.style   = {decoration={markings, mark=at position \anchopos with {
                      \draw[line width=0.6pt, -] (-0.1,-0.14) -- (0.1,0.14);
                      \node[font=\scriptsize, above=1.5pt] {#1};}},
                    postaction={decorate}},
  % la misma marca en una flecha vertical: el número va a la derecha
  anchov/.style  = {decoration={markings, mark=at position \anchopos with {
                      \draw[line width=0.6pt, -] (-0.14,-0.1) -- (0.14,0.1);
                      \node[font=\scriptsize, right=3pt] {#1};}},
                    postaction={decorate}},
  ancho en/.code = {\def\anchopos{#1}},
  sen/.style     = {font=\scriptsize\ttfamily, inner sep=1.5pt},
  nodo/.style    = {font=\tiny\ttfamily, inner sep=1pt, black!75},
  pin/.style     = {font=\scriptsize, align=center, inner sep=1pt},
  grupo/.style   = {draw=black!55, dashed, line width=0.7pt, rounded corners=4pt},
  canal/.style   = {grupo},
  contenedor/.style = {draw=black!80, line width=1.0pt, rounded corners=6pt},
  % marca numerada: cada número es una fila de la tabla del documento
  marca/.style   = {circle, fill=black, text=white, font=\scriptsize\bfseries,
                    inner sep=0pt, minimum size=4.3mm},
  %
  % DECISIONES (docs/mca/vista_general_mca.md). Cada opción es un bloque, todas
  % entran a un mux (\muxdec, abajo) y el registro que elige entra por abajo.
  % la opción que trae el reset (el comportamiento histórico): borde doble
  defecto/.style = {double=white, double distance=0.8pt, line width=0.45pt},
  % rótulo de la decisión, el mismo Dn del panorama y del documento
  decision/.style = {circle, fill=black, text=white, font=\small\bfseries,
                     inner sep=0pt, minimum size=7mm},
  %
  % LÓGICA (docs/mca/figuras/contabilidad/ y los detalles de las decisiones)
  % registro: caja gris con el triángulo de reloj en el borde de abajo
  registro/.style = {draw=black, line width=0.6pt, fill=grisr, align=center,
                     font=\scriptsize, inner sep=2pt, minimum width=#1,
                     minimum height=0.85cm,
                     path picture={\draw[line width=0.5pt]
                       ([xshift=-1.1mm]path picture bounding box.south) --
                       ([yshift=1.6mm]path picture bounding box.south) --
                       ([xshift=1.1mm]path picture bounding box.south);}},
  registro/.default = 0.8cm,
  % condición de una cadena if / else-if: "sí" sale a la derecha, "no" abajo
  pregunta/.style = {draw=black, line width=0.5pt, rounded corners=7pt,
                     fill=white, font=\scriptsize\ttfamily, align=center,
                     inner sep=2.5pt, minimum height=0.62cm, text width=#1},
  pregunta/.default = 1.9cm,
  si/.style      = {font=\tiny, inner sep=1pt, above, pos=0.12},
  no/.style      = {font=\tiny, inner sep=1pt, left, pos=0.45},
  %
  titulo/.style  = {anchor=north west, font=\small, inner sep=4pt, align=left},
  rotulo/.style  = {anchor=north west, font=\scriptsize\itshape, black!70,
                    inner sep=2pt, align=left},
}

% submódulo: nombre en monoespaciada; lógica propia: nombre en romana
\newcommand{\blk}[2]{{\ttfamily\bfseries #1}\\[1pt]#2}
\newcommand{\etq}[2]{{\bfseries #1}\\[1pt]#2}
\newcommand{\punto}[1]{\fill (#1) circle (1.3pt);}

% Parámetros en el formato del datasheet: NOMBRE=valor en monoespaciada, una
% entrada por celda, hasta cuatro columnas. Una columna: \prm{A=1\\B=2};
% dos: \prm{A=1 & B=2\\C=3}.
\newcommand{\prm}[1]{{\ttfamily\fontsize{5.6}{6.6}\selectfont
  \begin{tabular}{@{}l@{\hspace{7pt}}l@{\hspace{7pt}}l@{\hspace{7pt}}l@{}}#1\end{tabular}}}

\def\anchopos{0.5}

% Mux de una decisión: trapecio con el lado ancho (las entradas) a la izquierda.
% #1 = nombre, #2 = centro, #3 = alto en cm. Deja tres coordenadas:
%   (#1-in)  centro del lado de entradas: las flechas llegan a (#1-in |- opcion.east)
%   (#1-out) la salida
%   (#1-sel) mitad del lado de abajo: por ahí entra el registro que elige
% El rótulo de adentro es opcional: \muxdec{m}{(4,0)}{3.6} y después
% \node[...] at (m) {...}.
\newcommand{\muxdec}[3]{%
  \path #2 coordinate (#1);
  \draw[line width=0.6pt, fill=grisr]
    ($(#1)+(-0.45,{#3/2})$) -- ($(#1)+(0.45,{#3/2-0.35})$) --
    ($(#1)+(0.45,{-#3/2+0.35})$) -- ($(#1)+(-0.45,{-#3/2})$) -- cycle;
  \coordinate (#1-in)  at ($(#1)+(-0.45,0)$);
  \coordinate (#1-out) at ($(#1)+(0.45,0)$);
  \coordinate (#1-sel) at ($(#1)+(0,{-#3/2+0.175})$);}

% El mismo mux rotado 90 grados, para un flujo de arriba hacia abajo: entradas
% por el lado ancho de arriba, salida abajo y el registro que elige por la
% izquierda. #1 = nombre, #2 = centro, #3 = ancho en cm. Deja (#1-in) centro del
% lado de entradas, (#1-out) la salida y (#1-sel) mitad del lado izquierdo.
\newcommand{\muxdecv}[3]{%
  \path #2 coordinate (#1);
  \draw[line width=0.6pt, fill=grisr]
    ($(#1)+({-#3/2},0.45)$) -- ($(#1)+({#3/2},0.45)$) --
    ($(#1)+({#3/2-0.35},-0.45)$) -- ($(#1)+({-#3/2+0.35},-0.45)$) -- cycle;
  \coordinate (#1-in)  at ($(#1)+(0,0.45)$);
  \coordinate (#1-out) at ($(#1)+(0,-0.45)$);
  \coordinate (#1-sel) at ($(#1)+({-#3/2+0.175},0)$);}

% Leyenda de colores + nota. #1 = posición, #2 = color/texto, #3 = nota.
\newcommand{\leyenda}[3]{%
  \begin{scope}[shift={(#1)}, every node/.style={font=\scriptsize}]
    \coordinate (t0) at (-0.5,0);
    \foreach \c/\t [count=\i] in {#2} {
      \pgfmathtruncatemacro\j{\i-1}
      \node[draw, fill=\c, minimum width=4mm, minimum height=3mm, inner sep=0pt,
            right=5mm of t\j.east, anchor=west] (l\i) {};
      \node[right=1.5mm of l\i] (t\i) {\t};
    }
    \node[anchor=north west, black!70, text width=19cm, align=left]
      at (-0.05,-0.3) {#3};
  \end{scope}}

% Lo que dicen los parámetros, para la nota de la leyenda
\newcommand{\notaparam}{\texttt{NOMBRE=valor}: parámetros con que el padre
  instancia el bloque (se omiten los que sólo pasan el del padre sin cambio).}

\tikzset{
  reg/.style   = {registro=#1, minimum height=0.6cm, font=\scriptsize\ttfamily},
  reg/.default = 0.8cm,
  regt/.style  = {registro=#1, minimum height=0.5cm, font=\tiny\ttfamily},
  regt/.default = 0.75cm,
  % una línea de retardo con una toma dinámica: celdas en el borde de abajo
  linea/.style = {draw=black, line width=0.6pt, fill=azul!35, minimum width=1.4cm,
                  minimum height=0.7cm, inner sep=1pt, font=\small,
                  path picture={\foreach \k in {1,...,7}
                    \draw[line width=0.4pt, black!55]
                      ($(path picture bounding box.south west)!\k/8!(path picture bounding box.south east)$)
                      -- ++(0,0.16);}},
  chico/.style = {font=\scriptsize\ttfamily, inner sep=1pt},
  rot/.style   = {font=\tiny, inner sep=0.5pt, black!80},
  red/.style   = {ctrl},
  fondo f/.style = {draw=black!45, dashed, line width=0.5pt, rounded corners=3pt,
                    fill=#1, inner sep=0pt},
  rotf/.style  = {font=\scriptsize\itshape, black!75, inner sep=2pt},
  etapa l/.style = {draw=black!30, line width=0.5pt, dash pattern=on 2pt off 2pt},
  ciclo/.style = {font=\scriptsize, inner sep=1.5pt, fill=white},
}
% un nombre de red que entra a #1: desde la izquierda, o desde arriba o abajo
% con el nombre a la izquierda de la punta
\newcommand{\redIzq}[2]{\draw[red] ($(#1)+(-0.4,0)$) node[left, chico, black!80] {#2} -- (#1);}
\newcommand{\redArr}[2]{\draw[red] ($(#1)+(0,0.28)$) node[left, chico, black!80] {#2} -- (#1);}
\newcommand{\redAba}[2]{\draw[red] ($(#1)+(0,-0.28)$) node[left, chico, black!80] {#2} -- (#1);}
% una constante que entra a #1 desde la izquierda o desde arriba
\newcommand{\cteIzq}[2]{\draw[flujo] ($(#1)+(-0.3,0)$) node[left, cte] {#2} -- (#1);}
\newcommand{\cteArr}[2]{\draw[flujo] ($(#1)+(0,0.32)$) node[above, cte] {#2} -- (#1);}
% una salida con nombre: flecha corta y el nombre
\newcommand{\nace}[2]{\draw[flujo] (#1) -- ++(0.4,0) node[right, chico] {#2};}
% los rótulos 1/0 de las entradas de un mux, a #2 del centro
\newcommand{\unocero}[2]{%
  \node[font=\tiny] at ($(#1-in)+(0.13,#2)$) {1};
  \node[font=\tiny] at ($(#1-in)+(0.13,-#2)$) {0};}
\newcommand{\cerouno}[2]{%
  \node[font=\tiny] at ($(#1-in)+(0.13,#2)$) {0};
  \node[font=\tiny] at ($(#1-in)+(0.13,-#2)$) {1};}
% la regla de muestras: #1 = y, #2 = lista de x de los registros, #3 = y de abajo
\newcommand{\etapas}[3]{%
  \foreach \x in {#2} \draw[etapa l] (\x,#1) -- (\x,#3);}

\begin{document}
\begin{tikzpicture}[fondo]

\def\xin{-0.2}    % donde arrancan las entradas

% ===================================================== fondos por función ===
% primero, para que todo lo demás quede encima
\node[fondo f=azul!10, fit={(0.2,2.25) (13.2,14.85)}] (ftrap) {};
\node[fondo f=azul!10, fit={(13.4,10.95) (15.45,14.85)}] (fwin) {};
\node[fondo f=verde!15, fit={(0.2,-2.05) (9.0,1.65)}] (fslot) {};
\node[fondo f=naranja!18, fit={(9.25,-2.05) (15.45,1.65)}] (fpub) {};

\node[rotf, anchor=north east, fill=azul!10, align=right] at (13.1,14.78)
  {\texttt{mca\_trap} $\times$\,3, uno por forma\\[1pt]
   \prm{XW=15 & DAW=10 & MW=20\\ TW=18 & LAT=8}};
\node[rotf, anchor=north east, align=right] at (fwin.south east)
  {\texttt{mca\_window}:\\ la ventana de\\ cada meseta,\\ una por forma};
\node[rotf, anchor=north west] at (fslot.south west)
  {\texttt{mca\_slot} $\times$\,3: la suma sobre cada meseta, una por forma};
\node[rotf, anchor=north east] at (fpub.south east)
  {\texttt{mca\_publish}: el promedio de la forma elegida};

% ======================================================= la regla de muestras
\def\yra{15.15}
\etapas{15.0}{6.4, 9.8, 12.55}{5.6}
\node[ciclo] at (3.2,\yra) {muestra $n$};
\node[ciclo] at (8.1,\yra) {$n{+}1$};
\node[ciclo] at (11.2,\yra) {$n{+}2$};
\def\yrb{4.85}
\etapas{4.7}{2.35, 6.2, 8.15, 10.45, 12.55}{2.35}
\node[ciclo] at (1.0,\yrb) {$n{+}3$};
\node[ciclo] at (4.3,\yrb) {$n{+}4$};
\node[ciclo] at (7.2,\yrb) {$n{+}5$};
\node[ciclo] at (9.3,\yrb) {$n{+}6$};
\node[ciclo] at (11.5,\yrb) {$n{+}7$};
\node[ciclo] at (12.95,\yrb) {$n{+}8$};

% ================================================================ entradas ===
\foreach \y/\p/\n/\w in {14.2/cfg\_k\_i/k/10, 13.65/cfg\_l\_i/l/10,
                         13.1/cfg\_mq\_i/Mq/20, 12.55/cfg\_sh\_i/sh/6}
  \draw[flujo, ancho=\w, ancho en=0.45] (\xin,\y) node[left, sen] {\p} -- (0.75,\y)
    node[right, chico] {\n};
\draw[flujo] (\xin,12.0) node[left, sen] {clr\_i} -- (0.75,12.0) node[right, chico] {clr};
\draw[flujo] (\xin,11.45) node[left, sen] {val\_i} -- (0.75,11.45) node[right, chico] {val};

% ===================================== fila 1: las banderas de llenado ======
% cnt cuenta las muestras desde el clr; v[n-j] existe si cnt >= j
\node[oper] (acn) at (2.2,13.3) {$+$};
\node[satur] (scn) at (3.0,13.3) {};
\node[reg=0.8cm] (cnt) at (3.95,13.3) {cnt};
\node[oper] (gk)  at (5.35,14.3) {$\geq$};
\node[oper] (gl)  at (5.35,13.3) {$\geq$};
\node[oper] (gkl) at (5.35,12.3) {$\geq$};
\node[regt] (fk1)  at (6.4,14.3) {fk1};
\node[regt] (fl1)  at (6.4,13.3) {fl1};
\node[regt] (fkl1) at (6.4,12.3) {fkl1};
\node[regt] (fkl2) at (9.8,12.3) {fkl2};
\node[oper] (akl) at (2.2,12.15) {$+$};
\node[reg=0.75cm] (kl) at (3.25,12.15) {kl};

\cteArr{acn.north}{1}
\draw[flujo] (acn.east) -- (scn.west);
\draw[flujo] (scn.east) -- (cnt.west);
\draw[red] (3.75,13.88) node[left, chico, black!80] {val} -- (3.75,13.6);
\draw[red] (4.15,13.88) node[right, chico, black!80] {clr} -- (4.15,13.6);
\draw[line width=0.8pt] (cnt.east) -- (4.6,13.3);
\draw[line width=0.8pt] (4.6,14.3) -- (4.6,12.3);
\punto{4.6,13.3}
\draw[flujo] (4.6,14.3) -- (gk.west);
\draw[flujo] (4.6,13.3) -- (gl.west);
\draw[flujo] (4.6,12.3) -- (gkl.west);
\punto{4.6,12.85}
\draw[flujo] (4.6,12.85) -- (2.2,12.85) -- (acn.south);
\redArr{gk.north}{k}
\redArr{gl.north}{l}
\redArr{gkl.north}{kl}
\draw[flujo] (gk.east) -- (fk1.west);
\draw[flujo] (gl.east) -- (fl1.west);
\draw[flujo] (gkl.east) -- (fkl1.west);
\draw[flujo] (fkl1.east) -- (fkl2.west);
\nace{fk1.east}{fk1}
\nace{fl1.east}{fl1}
\nace{fkl2.east}{fkl2}
\redIzq{akl.west}{k}
\redAba{akl.south}{l}
\draw[flujo] (akl.east) -- (kl.west);
\nace{kl.east}{kl}

% ============================================ fila 1: las líneas de retardo ==
\node[reg] (v1) at (6.4,10.45) {v1};
\node[linea] (lk) at (2.9,9.35) {$z^{-k}$};
\node[reg] (rk) at (6.4,9.35) {rd\_k};
\node[linea] (ll) at (2.9,8.1) {$z^{-l}$};
\node[reg] (rl) at (6.4,8.1) {rd\_l};
\node[linea] (lkl) at (8.2,6.7) {$z^{-k}$};
\node[reg] (rkl) at (9.8,6.7) {rd\_kl};

\draw[flujo, ancho=15, ancho en=0.12] (\xin,10.45) node[left, sen] {x\_i} -- (v1.west);
\punto{1.25,10.45}
\draw[flujo] (1.25,10.45) |- (lk.west);
\punto{1.25,9.35}
\draw[flujo] (1.25,9.35) |- (ll.west);
\draw[flujo] (lk.east) -- (rk.west);
\draw[flujo] (ll.east) -- (rl.west);
\punto{6.95,8.1}
\draw[flujo] (6.95,8.1) |- (lkl.west);
\draw[flujo] (lkl.east) -- (rkl.west);
% cada línea escribe en wp y lee en wp - k (o wp - l)
\foreach \l/\a in {lk/a\_k, ll/a\_l, lkl/a\_k} {
  \draw[red] ($(\l.south west)+(0.3,-0.28)$) node[left, chico, black!80] {wp} -- ++(0,0.28);
  \draw[red] ($(\l.south east)+(-0.25,-0.28)$) node[left, chico, black!80] {\a} -- ++(0,0.28);
}

% el puntero y las dos direcciones de lectura
\node[oper] (awp) at (1.8,6.75) {$+$};
\node[reg=0.8cm] (wp) at (2.8,6.75) {wp};
\node[oper] (oak) at (4.35,6.75) {$-$};
\node[oper] (oal) at (4.35,6.0) {$-$};
\cteIzq{awp.west}{1}
\draw[flujo] (awp.east) -- (wp.west);
\draw[red] (2.6,6.17) node[left, chico, black!80] {val} -- (2.6,6.45);
\draw[red] (3.0,6.17) node[right, chico, black!80] {clr} -- (3.0,6.45);
\punto{3.55,6.75}
\draw[flujo] (wp.east) -- (oak.west);
\draw[flujo] (3.55,6.75) -- (3.55,7.2) -- (1.8,7.2) -- (awp.north);
\draw[flujo] (3.55,6.75) |- (oal.west);
\redArr{oak.north}{k}
\redAba{oal.south}{l}
\nace{oak.east}{a\_k}
\nace{oal.east}{a\_l}
\foreach \o in {oak, oal} {
  \node[rot, above left] at (\o.west) {$+$};
}
\node[rot, right] at ($(oak.north)+(0.04,0.1)$) {$-$};
\node[rot, right] at ($(oal.south)+(0.04,-0.1)$) {$-$};

% ======================================= fila 1: las tomas y la doble diferencia
% una toma que todavía no existe desde el clr lee 0
\muxdec{mk}{(7.65,9.15)}{0.9}
\unocero{mk}{0.2}
\muxdec{ml}{(8.05,7.9)}{0.9}
\unocero{ml}{0.2}
\muxdec{mkl}{(11.1,6.5)}{0.9}
\unocero{mkl}{0.2}
\node[oper] (re) at (8.75,10.45) {$-$};
\node[reg] (e2) at (9.8,10.45) {e2};
\node[reg] (l2) at (9.8,7.9) {l2};
\node[oper] (sd) at (11.6,7.9) {$+$};
\node[reg] (d3) at (12.55,7.9) {d3};

\draw[flujo] (rk.east) -- (mk-in |- 0,9.35);
\cteIzq{$(mk-in)+(0,-0.2)$}{0}
\draw[red] ($(mk)+(0,0.53)$) node[left, chico, black!80] {fk1} -- ($(mk)+(0,0.275)$);
\draw[flujo] (mk-out) -| (re.south);
\draw[flujo] (v1.east) -- (re.west);
\node[rot, above left] at (re.west) {$+$};
\node[rot, right] at ($(re.south)+(0.04,-0.1)$) {$-$};
\draw[flujo, ancho=16, ancho en=0.5] (re.east) -- (e2.west);

\draw[flujo] (rl.east) -- (ml-in |- 0,8.1);
\cteIzq{$(ml-in)+(0,-0.2)$}{0}
\redAba{$(ml)+(0,-0.275)$}{fl1}
\draw[flujo] (ml-out) -- (l2.west);

\draw[flujo] (rkl.east) -- (mkl-in |- 0,6.7);
\cteIzq{$(mkl-in)+(0,-0.2)$}{0}
\draw[red] ($(mkl)+(0,0.53)$) node[left, chico, black!80] {fkl2} -- ($(mkl)+(0,0.275)$);

% d = e2 - l2 + v[n-k-l]: una sola suma de tres
\draw[flujo] (e2.east) -| (sd.north);
\draw[flujo] (l2.east) -- (sd.west);
\draw[flujo] (mkl-out) -| (sd.south);
\node[rot, right] at ($(sd.north)+(0.04,0.12)$) {$+$};
\node[rot, above left] at (sd.west) {$-$};
\node[rot, right] at ($(sd.south)+(0.04,-0.12)$) {$+$};
\draw[flujo] (sd.east) -- (d3.west);

% d baja a la fila 2 por el canal de la derecha
\draw[flujo, ancho=17, ancho en=0.55] (d3.east) -- (13.0,7.9) -- (13.0,5.2) -- (0.45,5.2) -- (0.45,3.75) -- (0.99,3.75);
\node[chico, anchor=south] at (12.55,5.2) {d};

% ============================== fila 2: integradores, polo-cero y salida ====
\def\yb{3.75}
\node[oper] (ap) at (1.25,\yb) {$+$};
\node[reg] (p4) at (2.35,\yb) {p4};
\node[reg] (d4) at (2.35,2.65) {d4};
\node[desp] (dmf) at (3.95,\yb) {\texttt{<{}<{}< 4}};
\node[oper] (mul) at (3.95,2.65) {$\times$};
\node[oper] (ar) at (5.2,\yb) {$+$};
\node[reg] (r5) at (6.2,\yb) {r5};
\node[oper] (as) at (7.2,\yb) {$+$};
\node[reg] (s6) at (8.15,\yb) {s6};
\node[desp] (dsh) at (9.35,\yb) {\texttt{>{}>{}>}};
\node[reg] (q7) at (10.45,\yb) {q7};
\node[satur] (st) at (11.5,\yb) {};
\node[rot, below] at (st.south) {18 b};
\node[reg] (t8) at (12.55,\yb) {t8};

% el primer integrador: p <= p + d; d4 lo acompaña
\punto{0.75,\yb}
\draw[flujo] (0.75,\yb) |- (d4.west);
\draw[flujo] (ap.east) -- (p4.west);
\punto{3.05,\yb}
\draw[flujo, ancho=26, ancho en=0.68] (p4.east) -- (dmf.west);
\draw[flujo] (3.05,\yb) -- (3.05,3.25) -- (1.25,3.25) -- (ap.south);
\draw[red] ($(p4.north)+(0,0.3)$) node[above, chico, black!80] {clr} -- (p4.north);

% el polo-cero: r = p * 2^4 + Mq * d
\draw[flujo] (d4.east) -- (mul.west);
\redArr{mul.north}{Mq}
\draw[flujo] (dmf.east) -- (ar.west);
\draw[flujo] (mul.east) -| (ar.south);
\draw[flujo, ancho=39, ancho en=0.5] (ar.east) -- (r5.west);

% el segundo integrador: s <= s + r
\draw[flujo] (r5.east) -- (as.west);
\draw[flujo] (as.east) -- (s6.west);
\punto{8.8,\yb}
\draw[flujo, ancho=48, ancho en=0.55] (s6.east) -- (dsh.west);
\draw[flujo] (8.8,\yb) -- (8.8,3.25) -- (7.2,3.25) -- (as.south);
\draw[red] ($(s6.north)+(0,0.3)$) node[above, chico, black!80] {clr} -- (s6.north);

% corrimiento y saturación
\redArr{dsh.north}{sh}
\draw[flujo] (dsh.east) -- (q7.west);
\draw[flujo] (q7.east) -- (st.west);
\draw[flujo] (st.east) -- (t8.west);

% =================================== mca_window: la ventana de cada meseta ===
% lw a la izquierda (fin), off a la derecha (ini)
\node[oper] (aof) at (14.95,14.0) {$+$};
\node[satur] (sof) at (14.95,12.85) {};
\node[regt=0.95cm] (rini) at (14.95,11.85) {trap\_ini};
\node[desp] (dlw) at (13.9,14.45) {\texttt{1 <{}<}};
\node[oper] (afi) at (13.9,13.55) {$+$};
\node[satur] (sfi) at (13.9,12.85) {};
\node[regt=0.95cm] (rfin) at (13.9,11.85) {trap\_fin};

\draw[flujo, anchov=16, ancho en=0.55] (14.95,15.75) node[above left, sen, xshift=4pt] {cfg\_trap\_off\_i} -- (aof.north);
\cteIzq{aof.west}{8}
\punto{14.95,13.55}
\draw[flujo] (aof.south) -- (sof.north);
\draw[flujo] (14.95,13.55) -- (afi.east);
\draw[flujo] (sof.south) -- (rini.north);
\draw[flujo, anchov=3, ancho en=0.5] (13.9,15.25) node[above left, sen, xshift=4pt] {cfg\_trap\_lw\_i} -- (dlw.north);
\draw[flujo] (dlw.south) -- (afi.north);
\draw[flujo] (afi.south) -- (sfi.north);
\draw[flujo] (sfi.south) -- (rfin.north);
\draw[flujo] (rini.south) -- ++(0,-0.35) node[below, chico] {ini};
\draw[flujo] (rfin.south) -- ++(0,-0.35) node[below, chico] {fin};

% ======================================= mca_slot: la suma sobre la meseta ===
% en = acumula y ini <= len < fin
\node[oper] (ge) at (1.4,0.0) {$\geq$};
\node[oper] (lt) at (1.4,-0.9) {$<$};
\node[puerta y=nnn, anchor=input 1, minimum height=0.7cm] (yen) at (2.55,0.0) {};
\redIzq{ge.west}{len}
\redIzq{lt.west}{len}
\redArr{ge.north}{ini}
\redAba{lt.south}{fin}
\draw[flujo] (ge.east) -- (yen.input 1);
\draw[flujo] (lt.east) -- ++(0.45,0) |- (yen.input 2);
\draw[red] (2.4,-1.3) node[below, chico, black!80] {acumula} |- (yen.input 3);
\nace{yen.output}{en}

% la suma: al abrir, 0; con en, + el trapecio; si no, la que tenía
\node[oper] (asu) at (4.3,0.4) {$+$};
\muxdec{mse}{(5.4,0.6)}{0.9}
\cerouno{mse}{0.2}
\muxdec{msa}{(6.55,0.8)}{0.9}
\unocero{msa}{0.2}
\node[reg=0.9cm] (rsu) at (7.75,0.8) {suma};

% el trapecio llega de t8
\draw[flujo, ancho=18, ancho en=0.08] (t8.east) -- (13.1,\yb) -- (13.1,1.95) -- (3.6,1.95) -- (3.6,0.4) -- (asu.west);
\node[sen, right] at (13.1,2.9) {trap\_o};
\draw[flujo] (asu.east) -- (mse-in |- 0,0.4);
\draw[flujo] (mse-out) -- (msa-in |- 0,0.6);
\cteIzq{$(msa-in)+(0,0.2)$}{0}
\redAba{$(mse)+(0,-0.275)$}{en}
\redAba{$(msa)+(0,-0.275)$}{abre}
\draw[flujo] (msa-out) -- (rsu.west);
\punto{8.5,0.8}
\draw[flujo] (8.5,0.8) -- (8.5,1.45) -- (4.3,1.45) -- (asu.north);
\punto{4.75,1.45}
\draw[flujo] (4.75,1.45) |- (mse-in |- 0,0.8);

% la cuenta, con la misma regla
\node[oper] (acu) at (4.3,-1.45) {$+$};
\muxdec{mce}{(5.4,-1.25)}{0.9}
\cerouno{mce}{0.2}
\muxdec{mca}{(6.55,-1.05)}{0.9}
\unocero{mca}{0.2}
\node[reg=0.9cm] (rcu) at (7.75,-1.05) {cuenta};
\cteIzq{acu.west}{1}
\draw[flujo] (acu.east) -- (mce-in |- 0,-1.45);
\draw[flujo] (mce-out) -- (mca-in |- 0,-1.25);
\cteIzq{$(mca-in)+(0,0.2)$}{0}
\redAba{$(mce)+(0,-0.275)$}{en}
\redAba{$(mca)+(0,-0.275)$}{abre}
\draw[flujo] (mca-out) -- (rcu.west);
\punto{8.5,-1.05}
\draw[flujo] (8.5,-1.05) -- (8.5,-0.45) -- (4.3,-0.45) -- (acu.north);
\punto{4.75,-0.45}
\draw[flujo] (4.75,-0.45) |- (mce-in |- 0,-1.05);

% ========================================== mca_publish: el promedio ========
\muxdec{mfs}{(10.2,0.8)}{0.9}
\muxdec{mfn}{(10.2,-1.05)}{0.9}
\node[desp] (dpr) at (11.35,0.8) {\texttt{>{}>{}>}};
\node[satur] (spr) at (12.4,0.8) {};
\node[rot, above] at (spr.north) {$0\ldots2^{16}{-}1$};
\muxdec{mo}{(13.35,0.6)}{0.9}
\unocero{mo}{0.2}
\node[reg=1.0cm] (rft) at (14.4,0.6) {F\_TRAP};
\node[desp] (d2l) at (11.35,-0.25) {\texttt{1 <{}<}};
\node[oper] (eq) at (11.35,-1.05) {$=$};
\node[puerta y=ni, anchor=input 1] (yok) at (12.2,-1.05) {};

\draw[flujo, ancho=3$\times$25, ancho en=0.45] (8.5,0.8) -- (mfs-in);
\draw[flujo, ancho=3$\times$8, ancho en=0.45] (rcu.east) -- (mfn-in);
\redAba{$(mfs)+(0,-0.275)$}{f*}
\redAba{$(mfn)+(0,-0.275)$}{f*}
\draw[flujo] (mfs-out) -- (dpr.west);
\redArr{dpr.north}{lw}
\draw[flujo] (dpr.east) -- (spr.west);
\draw[flujo] (spr.east) -- (mo-in |- 0,0.8);
\cteIzq{$(mo-in)+(0,-0.2)$}{0}
\draw[flujo] (mfn-out) -- (eq.west);
\redIzq{d2l.west}{lw}
\draw[flujo] (d2l.south) -- (eq.north);
\draw[flujo] (eq.east) -- (yok.input 1);
\draw[red] (11.95,-1.75) node[below, chico, black!80] {marcado} |- (yok.input 2);
\draw[flujo] (yok.output) -| ($(mo)+(0,-0.275)$);
\draw[flujo] (mo-out) -- (rft.west);
\draw[flujo, ancho=16, ancho en=0.45] (rft.east) -- (15.3,0.6);

% ================================================================= leyenda ===
\leyenda{0.2,-2.95}{azul!10/preprocesamiento, verde!15/estadísticos del pulso,
  naranja!18/publicación, grisr/registro}{}

\end{tikzpicture}
\end{document}
